\documentclass[12pt,a4paper]{article}

\usepackage[utf8]{inputenc}
\usepackage[vietnamese]{babel}
\usepackage{amsmath,amssymb,amsthm}
\usepackage{geometry}
\usepackage{enumitem}
\usepackage[most]{tcolorbox}
\usepackage{hyperref}
\usepackage{fancyhdr}
\usepackage{tikz}
\usepackage{eso-pic}
\usepackage{xcolor}

\geometry{a4paper, margin=2cm, bottom=2.5cm}
\hypersetup{colorlinks=true, linkcolor=blue!70!black}
\setlist[itemize]{topsep=4pt, itemsep=2pt}
\setlist[enumerate]{topsep=4pt, itemsep=2pt}

\AddToShipoutPictureFG{%
  \AtPageCenter{%
    \begin{tikzpicture}[remember picture, overlay]
      \node[rotate=45, scale=1, opacity=0.06]{\fontsize{60}{72}\selectfont\bfseries\sffamily Đặng Tấn Quí};
    \end{tikzpicture}%
  }%
}

\pagestyle{fancy}
\fancyhf{}
\fancyhead[L]{\small\textit{Toán cao cấp 2 -- Giải tích một biến}}
\fancyhead[R]{\small\textit{Đặng Tấn Quí}}
\fancyfoot[C]{\thepage}
\fancyfoot[L]{\footnotesize\textcopyright\ Đặng Tấn Quí}
\fancyfoot[R]{\footnotesize SĐT: 0377\,122\,210}
\renewcommand{\headrulewidth}{0.4pt}
\renewcommand{\footrulewidth}{0.4pt}

\fancypagestyle{plain}{\fancyhf{}\fancyfoot[C]{\thepage}\fancyfoot[L]{\footnotesize\textcopyright\ Đặng Tấn Quí}\fancyfoot[R]{\footnotesize SĐT: 0377\,122\,210}\renewcommand{\headrulewidth}{0pt}\renewcommand{\footrulewidth}{0.4pt}}

\newtcolorbox{dinhnghia}[1][]{enhanced, breakable, colback=blue!3!white, colframe=blue!60!black, fonttitle=\bfseries, title={Định nghĩa}, #1}
\newtcolorbox{dinhly}[1][]{enhanced, breakable, colback=green!3!white, colframe=green!50!black, fonttitle=\bfseries, title={Định lý}, #1}
\newtcolorbox{tinhchat}[1][]{enhanced, breakable, colback=yellow!5!white, colframe=orange!50!black, fonttitle=\bfseries, title={Tính chất}, #1}
\newtcolorbox{phuongphap}[1][]{enhanced, breakable, colback=gray!5!white, colframe=gray!50!black, fonttitle=\bfseries, title={Phương pháp}, #1}
\newtcolorbox{luuy}[1][]{enhanced, breakable, colback=red!3!white, colframe=red!50!black, fonttitle=\bfseries, title={Lưu ý}, #1}
\newtcolorbox{congthuc}[1][]{enhanced, breakable, colback=cyan!3!white, colframe=cyan!50!black, fonttitle=\bfseries, title={Công thức}, #1}
\newtcolorbox{vidu}[1][]{enhanced, breakable, colback=violet!3!white, colframe=violet!50!black, fonttitle=\bfseries, title={Ví dụ}, #1}

\begin{document}

\thispagestyle{empty}
\begin{tikzpicture}[remember picture, overlay]
  \fill[blue!80!black] (current page.south west) rectangle (current page.north east);
  \node[anchor=west, white, font=\fontsize{26}{32}\bfseries\selectfont]
    at ([xshift=3cm, yshift=-7.5cm]current page.north west) {TOÁN CAO CẤP 2};
  \node[anchor=west, white, font=\fontsize{18}{24}\selectfont]
    at ([xshift=3cm, yshift=-9cm]current page.north west) {Giải tích một biến --- Tích phân, chuỗi --- Năm 1, HK2};
  \node[anchor=west, white, font=\Large\bfseries] at ([xshift=3cm, yshift=-13cm]current page.north west) {Biên soạn:};
  \node[anchor=west, yellow!80!white, font=\fontsize{22}{28}\bfseries\selectfont] at ([xshift=3cm, yshift=-14.5cm]current page.north west) {ĐẶNG TẤN QUÍ};
  \node[anchor=south east, white, font=\Large\bfseries] at ([xshift=-2cm, yshift=3cm]current page.south east) {\the\year};
\end{tikzpicture}

\newpage
\tableofcontents
\newpage

\section{Tích phân}

\subsection{Tích phân bất định}

\begin{dinhnghia}
  $\displaystyle\int f(x)\,dx = F(x) + C$ với $F'(x) = f(x)$.
\end{dinhnghia}

\begin{phuongphap}[title={Đổi biến}]
  Đặt $u = g(x)$ $\Rightarrow$ $du = g'(x)dx$. $\displaystyle\int f(g(x))g'(x)\,dx = \int f(u)\,du$.
\end{phuongphap}

\begin{phuongphap}[title={Tích phân từng phần}]
  $\displaystyle\int u\,dv = uv - \int v\,du$. Thường dùng khi có $\ln x$, $e^x$, $\sin x$, $\cos x$.
\end{phuongphap}

\subsection{Tích phân xác định}

\begin{congthuc}
  $\displaystyle\int_a^b f(x)\,dx = F(b) - F(a)$, $\quad \int_a^b f = -\int_b^a f$, $\quad \int_a^c f = \int_a^b f + \int_b^c f$.
\end{congthuc}

\begin{dinhly}[title={Đổi biến trong tích phân xác định}]
  Đặt $x = \varphi(t)$, $a = \varphi(\alpha)$, $b = \varphi(\beta)$:
  \[
    \int_a^b f(x)\,dx = \int_\alpha^\beta f(\varphi(t))\,\varphi'(t)\,dt.
  \]
\end{dinhly}

\subsection{Tích phân suy rộng}

\begin{dinhnghia}
  \textbf{Loại 1} (miền vô hạn): $\displaystyle\int_a^{+\infty} f(x)\,dx = \lim_{b \to +\infty} \int_a^b f(x)\,dx$.
\end{dinhnghia}

\begin{dinhnghia}
  \textbf{Loại 2} (hàm không bị chặn): Nếu $f$ không bị chặn tại $b$: $\displaystyle\int_a^b f(x)\,dx = \lim_{t \to b^-} \int_a^t f(x)\,dx$.
\end{dinhnghia}

\begin{dinhly}[title={Tiêu chuẩn hội tụ}]
  \begin{itemize}
    \item $\displaystyle\int_1^{+\infty} \frac{1}{x^p}\,dx$ hội tụ khi $p > 1$, phân kỳ khi $p \le 1$.
    \item So sánh: Nếu $0 \le f(x) \le g(x)$ và $\int g$ hội tụ thì $\int f$ hội tụ.
  \end{itemize}
\end{dinhly}

\section{Chuỗi số}

\subsection{Chuỗi số thực}

\begin{dinhnghia}
  Chuỗi $\displaystyle\sum_{n=1}^{\infty} a_n$ hội tụ $\Leftrightarrow$ dãy tổng riêng $S_n = \sum_{k=1}^n a_k$ có giới hạn hữu hạn.
\end{dinhnghia}

\begin{dinhly}[title={Điều kiện cần}]
  $\sum a_n$ hội tụ $\Rightarrow$ $\lim_{n \to \infty} a_n = 0$.
\end{dinhly}

\begin{dinhly}[title={Chuỗi dương}]
  \begin{itemize}
    \item So sánh: $0 \le a_n \le b_n$, $\sum b_n$ hội tụ $\Rightarrow$ $\sum a_n$ hội tụ.
    \item D'Alembert: $\displaystyle L = \lim \frac{a_{n+1}}{a_n}$: $L < 1$ hội tụ, $L > 1$ phân kỳ.
    \item Cauchy: $\displaystyle L = \lim \sqrt[n]{a_n}$: $L < 1$ hội tụ, $L > 1$ phân kỳ.
    \item Tích phân: Nếu $a_n = f(n)$ với $f$ dương, giảm, liên tục thì $\sum a_n$ và $\int_1^{+\infty} f(x)\,dx$ cùng hội tụ hoặc cùng phân kỳ.
  \end{itemize}
\end{dinhly}

\begin{congthuc}
  Chuỗi $p$: $\displaystyle\sum_{n=1}^{\infty} \frac{1}{n^p}$ hội tụ $\Leftrightarrow$ $p > 1$.
\end{congthuc}

\subsection{Chuỗi hàm}

\begin{dinhnghia}
  \textbf{Chuỗi lũy thừa}: $\displaystyle\sum_{n=0}^{\infty} a_n (x - x_0)^n$. Bán kính hội tụ $R$: hội tụ khi $|x - x_0| < R$, phân kỳ khi $|x - x_0| > R$.
\end{dinhnghia}

\begin{congthuc}[title={Công thức Cauchy--Hadamard}]
  $R = \dfrac{1}{\displaystyle\limsup_{n \to \infty} \sqrt[n]{|a_n|}}$, hoặc $R = \displaystyle\lim_{n \to \infty} \left|\frac{a_n}{a_{n+1}}\right|$ nếu giới hạn tồn tại.
\end{congthuc}

\begin{vidu}
  $\displaystyle e^x = \sum_{n=0}^{\infty} \frac{x^n}{n!}$, $\quad \sin x = \sum_{n=0}^{\infty} \frac{(-1)^n x^{2n+1}}{(2n+1)!}$, $\quad \frac{1}{1-x} = \sum_{n=0}^{\infty} x^n$ ($|x| < 1$).
\end{vidu}

\section{Phương trình vi phân (mở đầu)}

\begin{dinhnghia}
  \textbf{Phương trình vi phân cấp 1}: $F(x, y, y') = 0$. Dạng $y' = f(x,y)$ gọi là \textbf{dạng chuẩn}.
\end{dinhnghia}

\begin{phuongphap}[title={Phương trình tách biến}]
  $y' = g(x)h(y)$ $\Rightarrow$ $\displaystyle\int \frac{dy}{h(y)} = \int g(x)\,dx$.
\end{phuongphap}

\begin{phuongphap}[title={Phương trình tuyến tính cấp 1}]
  $y' + p(x)y = q(x)$. Nghiệm: $y = e^{-\int p}\left(\int q\,e^{\int p}\,dx + C\right)$.
\end{phuongphap}

\end{document}
